Detector bias and what a calibration inherits from it, on rendered images with known camera and poses. Hue encodes the direction of a vector in the image and saturation its length; the disc is the key, the number the length at full saturation (\emph{Geometric} and \emph{Radiometric} share the key between them). The fields are means per sensor bin, interpolated bilinearly. \emph{True Projected vs.\ Detected}: detected minus true position, the mean over all detections of a bin (OpenCV + AprilGrid: the tag corners; Sch\"ops + Siemens stars: the star centres). \emph{Calibrated Camera \& Pose}: board points projected with the calibrated camera and poses, minus their true pixels. \emph{Calibrated Camera}: the same under the true poses, one global rotation removed. \emph{Geometric}: the detector-based calibration; \emph{Radiometric}: after our radiometric refinement. A key with a band around its disc and $\geq$ (Sch\"ops + Siemens stars) is cut: its radius is the 98\,\% quantile of the bin vectors, and the longer ones (5 and 5 bins, up to $0.14$\,px) are drawn in the full colour of their direction.
